\section{卷积神经网络（CNN）架构}

\subsection{CNN 的三个核心组件}

了解了卷积操作后，我们可以构建完整的卷积神经网络。一个典型的 CNN 由三种层交替组成：

\begin{figure}[H]
\centering
\includegraphics[width=0.9\textwidth]{slides-images/slide-44.jpg}
\caption{CNN 分类网络的完整流程：输入图像 $\rightarrow$ 卷积层（生成特征图）$\rightarrow$ 最大池化 $\rightarrow$ $\cdots$ $\rightarrow$ 全连接层 $\rightarrow$ 分类输出。\protect\footnotemark}
\end{figure}
\footnotetext{Slides 第44页。}

\begin{importantbox}{CNN 的三个核心操作}
\begin{enumerate}
    \item \textbf{卷积}（Convolution）：应用滤波器生成特征图，提取局部特征
    \item \textbf{非线性激活}（Non-linearity）：通常使用 ReLU，引入非线性表达能力
    \item \textbf{池化}（Pooling）：对每个特征图进行下采样，降低维度并获得空间不变性
\end{enumerate}
通过训练，CNN 自动学习卷积层中滤波器的权重。
\end{importantbox}

\subsection{卷积层：局部连接}

卷积层中的每个神经元只与输入的一个局部区域相连。对于隐藏层中位置 $(p,q)$ 的神经元：

\begin{figure}[H]
\centering
\includegraphics[width=0.85\textwidth]{slides-images/slide-46.jpg}
\caption{卷积层的局部连接：每个神经元从局部区域取输入，计算加权和，加偏置。其数学表达与全连接层相同，但作用范围是局部的。\protect\footnotemark}
\end{figure}
\footnotetext{Slides 第46页。}

$$y_{p,q} = \sum_{i=1}^{k} \sum_{j=1}^{k} w_{ij} \cdot x_{i+p, j+q} + b$$

这个过程包括三步：(1) 应用权重窗口，(2) 计算线性组合，(3) 使用非线性函数激活。

\subsection{输出体积的空间排列}

每一层的输出不是一个二维平面，而是一个三维\textbf{体积}（volume），维度为 $h \times w \times d$：

\begin{figure}[H]
\centering
\includegraphics[width=0.9\textwidth]{slides-images/slide-47.jpg}
\caption{CNN 输出体积的维度：$h \times w \times d$，其中 $h$ 和 $w$ 是空间维度，$d$（深度）等于该层的滤波器数量。\protect\footnotemark}
\end{figure}
\footnotetext{Slides 第47页。}

\begin{itemize}
    \item $h$ 和 $w$：输出的空间维度（由输入尺寸、滤波器大小和步幅决定）
    \item $d$：深度，等于该层使用的\textbf{滤波器数量}
    \item \textbf{Stride}（步幅）：滤波器的滑动步长
    \item \textbf{Receptive Field}（感受野）：输入图像中与某个输出节点相连的区域
\end{itemize}

在 TensorFlow 中：\texttt{tf.keras.layers.Conv2D(filters=d, kernel\_size=(h,w), strides=s)}

在 PyTorch 中：\texttt{torch.nn.Conv2d(in\_channels, out\_channels=d, kernel\_size, stride=s)}

\begin{warningbox}{PyTorch 需要指定输入通道数}
TensorFlow/Keras 的 \texttt{Conv2D} 会自动推断输入维度，但 PyTorch 的 \texttt{Conv2d} 要求显式指定 \texttt{in\_channels}。初学者经常忘记这一点导致报错。
\end{warningbox}

\subsection{非线性激活函数：ReLU}

在每次卷积操作之后，需要应用非线性激活函数。CNN 中最常用的是\textbf{ReLU}（Rectified Linear Unit）：

\begin{figure}[H]
\centering
\includegraphics[width=0.85\textwidth]{slides-images/slide-48.jpg}
\caption{ReLU 激活函数：将所有负值替换为零，保留正值不变。这是一个逐像素的非线性操作。\protect\footnotemark}
\end{figure}
\footnotetext{Slides 第48页。}

$$g(z) = \max(0, z)$$

ReLU 的作用是将特征图中的所有负值替换为零，只保留正值（即检测到特征的区域）。为什么在图像任务中特别适合使用 ReLU？因为当滤波器检测到某个模式时输出为正值（匹配），而未检测到时输出为负值或零（不匹配），ReLU 恰好保留了"有检测到"的信号。

\subsection{池化：降维与空间不变性}

\textbf{池化}（Pooling）是卷积层之后的下采样操作。最常用的是\textbf{最大池化}（Max Pooling）：

\begin{figure}[H]
\centering
\includegraphics[width=0.85\textwidth]{slides-images/slide-49.jpg}
\caption{最大池化示例：使用 2$\times$2 窗口、步幅 2，在每个窗口中取最大值。4$\times$4 输入变为 2$\times$2 输出。\protect\footnotemark}
\end{figure}
\footnotetext{Slides 第49页。}

\begin{importantbox}{池化的两大作用}
\begin{enumerate}
    \item \textbf{降低维度}（Reduced dimensionality）：减小特征图的空间尺寸，从而减少后续层的计算量和参数量
    \item \textbf{空间不变性}（Spatial invariance）：由于取最大值，特征的精确位置变得不那么重要，增强了对位移的鲁棒性
\end{enumerate}
\end{importantbox}

Amini 提到，除了 Max Pooling，还可以问一个更一般的问题："还有什么其他方式可以实现下采样并保持空间不变性？"这暗示了后续课程可能介绍的其他技术，如带步幅的卷积（strided convolution）等。

\subsection{深层 CNN 中的表示学习}

通过堆叠多个"卷积 + ReLU + 池化"模块，CNN 能够学习越来越抽象的特征表示：

\begin{figure}[H]
\centering
\includegraphics[width=0.85\textwidth]{slides-images/slide-50.jpg}
\caption{深层 CNN 的表示学习：第一层学习边缘和暗斑等低层特征，第二层组合为眼睛、鼻子等中层特征，第三层组合为完整面部结构等高层特征。\protect\footnotemark}
\end{figure}
\footnotetext{Slides 第50页。参考 Lee+ ICML 2009。}

\begin{knowledgebox}{特征的层级组合}
深层 CNN 之所以强大，是因为每一层都在前一层检测到的特征基础上进行\textbf{组合}：
\begin{itemize}
    \item 第一个卷积层检测原始特征（如各种方向的边缘）
    \item 池化后，特征图缩小，相当于在更大的空间尺度上观察
    \item 第二个卷积层在缩小后的特征图上检测，相当于组合多个低层特征形成中层特征
    \item 如此重复，特征越来越抽象、语义越来越丰富
\end{itemize}
这种从简单到复杂的层级特征学习，正是深度学习中"深度"一词的含义。
\end{knowledgebox}

\subsection{本章小结}

一个完整的 CNN 分类网络由特征学习部分和分类部分组成。特征学习部分通过交替使用卷积层（提取局部特征）、非线性激活（引入表达力）和池化层（降维并增加不变性），逐层构建越来越抽象的特征表示。分类部分则将学习到的特征展平后通过全连接层输出类别概率。

\newpage
